%! Function Model Selection and Assumption Articulation — Unit 1, Lesson 1.7 — Warm-Up
\documentclass[10pt]{article}
\usepackage{precalculus-article}
\usepackage{precalculus-boxes}
\newcommand{\TallMath}[1]{$\displaystyle #1\rule[-1.4em]{0pt}{3.2em}$}

\begin{document}

\pageheader{Unit 1, Lesson 1.7}{Warm-Up}
\namedateperiod

\vspace{0.12in}

\begin{enumerate}[itemsep=14pt]

\item The table below shows equally-spaced inputs and their outputs.
      Complete the \textbf{first differences} column.

      \medskip
      \renewcommand{\arraystretch}{1.4}
      \begin{tabular}{|c|c|c|}
      \hline
      $x$ & $y$ & $\Delta y$ (1st diff.) \\
      \hline
      0 & 3  & ---         \\
      \hline
      1 & 7  & \blank{2cm} \\
      \hline
      2 & 11 & \blank{2cm} \\
      \hline
      3 & 15 & \blank{2cm} \\
      \hline
      4 & 19 & \blank{2cm} \\
      \hline
      \end{tabular}

      \medskip
      The first differences are \blank{2.2cm}. This means the rate of change is
      \blank{3cm}, so the data is best modeled by a \blank{3cm} function.

\item A second data set has first differences of $2, 4, 6, 8$.
      Are these first differences constant? \blank{1.5cm}

      \medskip
      Compute the \textbf{second differences} below:

      \medskip
      \renewcommand{\arraystretch}{1.4}
      \begin{tabular}{|c|c|c|}
      \hline
      1st diff. & $\Delta^2 y$ (2nd diff.) \\
      \hline
      2  & ---         \\
      \hline
      4  & \blank{2cm} \\
      \hline
      6  & \blank{2cm} \\
      \hline
      8  & \blank{2cm} \\
      \hline
      \end{tabular}

      \medskip
      The second differences are \blank{2.2cm}. This means the data is best modeled
      by a \blank{3cm} function.

\item A function's outputs change by roughly the same amount for each unit increase in input.
      What type of function is the best model? Circle one.

      \medskip
      \textbf{A.} Linear \qquad \textbf{B.} Quadratic \qquad \textbf{C.} Cubic \qquad \textbf{D.} Piecewise

\end{enumerate}

\end{document}
